\chapter{Guida Utente}

HearthCode è una versione semplificata di un gioco di carte a turni ispirato a
Hearthstone. Il giocatore affronta un avversario controllato dal computer,
giocando creature dalla propria mano, attaccando il nemico e cercando di
portare i suoi punti vita a zero.

\section{Navigazione nel menu}

\begin{figure}[htbp]
    \centering
    \includegraphics[width=0.75\textwidth]{images/Menu.png}
    \caption{Menu principale dell'applicazione.}
    \label{fig:guida-menu}
\end{figure}

All'avvio dell'applicazione viene mostrato il menu principale. Da qui è
possibile iniziare una nuova partita, consultare il mazzo oppure chiudere il
gioco. Il pulsante \textit{Play} porta alla scelta della difficoltà, mentre
\textit{Deck} apre una schermata riassuntiva con tutte le carte disponibili.

Nella schermata di selezione della difficoltà si può scegliere tra
\textit{Normal} e \textit{Hard}. La modalità \textit{Normal} usa un avversario
più semplice, mentre \textit{Hard} rende le decisioni del computer più
strategiche. Il pulsante \textit{Back} riporta al menu principale.

\begin{figure}[htbp]
    \centering
    \includegraphics[width=0.75\textwidth]{images/Difficulty.png}
    \caption{Schermata per la scelta della difficoltà.}
    \label{fig:guida-difficulty}
\end{figure}

La sezione \textit{Deck} permette di visualizzare le creature presenti nel
gioco. Ogni carta mostra costo in mana, nome, attacco e salute: queste
informazioni aiutano a capire quali creature conviene giocare durante la
partita. Anche da questa schermata il pulsante \textit{Back} consente di
tornare al menu.

\begin{figure}[htbp]
    \centering
    \includegraphics[width=0.75\textwidth]{images/Deck.png}
    \caption{Visualizzazione delle carte disponibili nel mazzo.}
    \label{fig:guida-deck}
\end{figure}

\section{Schermata di gioco}

\begin{figure}[htbp]
    \centering
    \includegraphics[width=0.75\textwidth]{images/Match.png}
    \caption{Schermata principale della partita.}
    \label{fig:guida-match}
\end{figure}

La partita è divisa in due aree: in alto si trovano le informazioni
dell'avversario, in basso quelle del giocatore. Per entrambi sono indicati i
punti vita, il mana disponibile, il numero di carte in mano, il numero di
creature sul campo e le carte rimaste nel mazzo. Le carte dell'avversario in
mano sono coperte, mentre quelle del giocatore sono visibili.

Al centro della schermata si trovano i campi di battaglia: \textit{Enemy Army}
per le creature nemiche e \textit{Player Army} per quelle del giocatore. A
sinistra sono presenti i pulsanti d'azione:
\begin{itemize}
    \item \textit{Place Card}: gioca sul campo la carta selezionata dalla mano,
    se il mana disponibile è sufficiente.
    \item \textit{Attack Hero}: attacca direttamente l'avversario con una
    creatura selezionata.
    \item \textit{Attack Creature}: attacca una creatura nemica dopo aver
    selezionato una propria creatura e un bersaglio avversario.
    \item \textit{End Turn}: termina il turno del giocatore e lascia agire il
    computer.
    \item \textit{Exit}: abbandona la partita e torna al menu, dopo una
    conferma.
\end{itemize}

\section{Come si gioca}

Ogni turno il giocatore dispone di una quantità di mana, indicata nella barra
in basso a sinistra. Per giocare una carta, bisogna selezionarla dalla mano e
premere \textit{Place Card}; la carta viene spostata nella propria armata e il
suo costo viene sottratto dal mana disponibile. Le carte appena posizionate non
possono attaccare immediatamente: diventano utilizzabili nei turni successivi.

Per attaccare, si seleziona una creatura già pronta nel proprio campo. Se si
vuole colpire direttamente l'avversario si preme \textit{Attack Hero}; se
invece si vuole colpire una creatura nemica, si seleziona anche il bersaglio
nell'\textit{Enemy Army} e poi si preme \textit{Attack Creature}. Dopo un
attacco la creatura viene esaurita e non può attaccare di nuovo nello stesso
turno.

Quando non si desidera compiere altre azioni, si preme \textit{End Turn}. Il
computer esegue il proprio turno automaticamente e poi il controllo torna al
giocatore. La partita termina quando i punti vita di uno dei due giocatori
raggiungono lo zero; viene quindi mostrata la schermata finale con il vincitore.
